\documentclass[conference]{IEEEtran}

\usepackage[utf8]{inputenc}
\usepackage[T1]{fontenc}
\usepackage{booktabs}
\usepackage{graphicx}
\usepackage{amsmath}
\usepackage{algorithm}
\usepackage{algpseudocode}

\begin{document}

\title{CCV: Games@Cloud Project Report}

\author{
  \IEEEauthorblockN{João Gouveia}
  \IEEEauthorblockA{
    Instituto Superior Técnico\\
    Lisbon, Portugal\\
    ist1102611@tecnico.ulisboa.pt
  }
  \and
  \IEEEauthorblockN{Francisco Fonseca}
  \IEEEauthorblockA{
    Instituto Superior Técnico\\
    Lisbon, Portugal\\
    ist1102492@tecnico.ulisboa.pt
  }
  \and
  \IEEEauthorblockN{Gonçalo Rua}
  \IEEEauthorblockA{
    Instituto Superior Técnico\\
    Lisbon, Portugal\\
    ist1102604@tecnico.ulisboa.pt
  }
}

\maketitle

\section{Introduction}
% TODO

\section{Architecture Overview}
% TODO

\section{Instrumentation}
% TODO

\subsection{Data Structures}
% TODO

\section{Workload And Request Cost Estimation}
% TODO

\section{Load Balancer}
% TODO

\section{Auto-Scaler}
We implemented a cost-aware auto-scaler based on a simplified version of
Kingfisher \cite{kingfisher}.

\subsection{Kingfisher}
Kingfisher uses integer linear programming (ILP) to efficiently determine the
optimal VM configuration for a given workload, taking into account either the
cost of the new configuration or the latency introduced by the transition.

Let $\lambda$ denote the number of requests an application is expected to handle
while meeting a target response time, and let $N$ be the maximum number of
servers that can be used to meet this goal. Let $M$ denote the number of
available VM types, and for a worker of type $j$, let $C_j$ denote its maximum
request rate and $p_j$ its cost. Let $T$ be the number of different transition
methods supported by the chosen cloud provider, and let $m_{ijk}$ denote the
latency introduced by transitioning server $i$ to type $j$ using transition
mechanism $k$.

The constraints are
\begin{align}
\sum_{i = 1}^N \sum_{j = 1}^M \sum_{k = 1}^T x_{ijk} C_j \ge; \lambda \
\end{align}
\begin{align}
\sum_{j = 1}^M \sum_{k = 1}^T x_{ijk} = 1,; \forall i
\end{align}

where $x_{ijk} \in \{0, 1\}$ is a variable indicating whether server $i$ is
transformed into a server of type $j$ using transition mechanism $k$.

The objective function is either
\begin{align}
\min \sum_{i = 1}^N \sum_{j = 1}^M \sum_{k = 1}^T p_j x_{ijk}
\end{align}
if we aim to minimize the configuration cost, or
\begin{align}
\min \sum_{i = 1}^N \sum_{j = 1}^M \sum_{k = 1}^T m_{ijk} x_{ijk}
\end{align}
if we aim to minimize transition latency.

\subsection{Our Solution}

According to the project guidelines, we have only one type of VM available
(\texttt{t2.micro}). While we are not required to support different transition
mechanisms, we must be able to minimize both startup latency and infrastructure
cost, rather than optimizing for just one.

As such, our encoding differs from Kingfisher's. We base our model on a simpler
formulation --- also presented in the paper --- that does not account for
varying transition latencies.

As previously mentioned, our load balancer is responsible for launching Lambda
functions, which handles startup latency. Consequently, our auto-scaler focuses
on minimizing the cost of the selected VM setup while ensuring that the system
can handle the workload at any given time.

Instead of targeting a specific request rate, as Kingfisher does, we opted to
use the system workload estimated by the load balancer. This decision was made
for the following reasons:
\begin{enumerate}
    \item We are already collecting metrics and estimating the workload on the
    load balancer, so no additional overhead is introduced.
    \item Using real-time information to guide scaling decisions is more
    accurate and flexible than determining a fixed limit prior to deployment.
\end{enumerate}

Let $\lambda$ denote the workload our application is subject to at a given
moment. For a VM of type $j$, let $C_j$ denote the maximum workload it can
handle and $p_j$ its cost.

Our constraint is:
\begin{align}
\sum_{j = 1}^M n_j C_j ;\ge; \lambda
\end{align}
where $n_j$ represents the number of workers of type $j$ in the new
configuration.

Our objective function is:
\begin{align}
\min \sum_{j = 1}^M n_j p_j
\end{align}

Our approach to auto-scaling is reactive: the algorithm is triggered
periodically, and—given the current estimated workload—it decides whether to
scale in, scale out, or maintain the current configuration.

\begin{thebibliography}{1}
\bibitem{kingfisher}
U. Sharma, P. Shenoy, S. Sahu, and A. Shaikh,
``Kingfisher: Cost-aware elasticity in the cloud'' in
\emph{Proc.\ IEEE INFOCOM}, Shanghai, China, Apr.\ 2011, pp.\ 206–210.
doi:10.1109/INFCOM.2011.5935016.

\end{thebibliography}

\end{document}
